\chapter{字形资源与质量评价}\label{ch:resources}
\section{本机大容量字体（0.1.118）}
\entry{主窗口 → 字体库 → 制作字体}
预计候选超过 30,000 份时，制作前会提示。选择\ui{保留全部，继续制作}可保留原定类型和全部候选。候选较少但含大量来源详情时，也可主动勾选\ui{使用本机大容量模式}。导入本机大容量包同样需要确认。大型字体的制作、校验和导入可能占用较多内存、磁盘空间与时间。
字体库会标明\ui{本机大容量}。需要跨软件使用时，从\ui{导出字体包}选择\ui{通用 v1 副本（选择候选）}。可以手动勾选，或先为每字每种字形各选一份，再逐项调整。导出的副本依通用规格校验。

\section{生成后的评分}
在生成页勾选来源分、多样性分或两者。「计算方式」也控制评分进程。自动模式在大批次使用最多 8 个进程，实际数量受处理器与可用内存限制。评分阶段单独显示已评分字数，取消后仍可使用已生成的字形。

\section{导入统一字库（0.1.117）}
\entry{主窗口 → 字体库 → 导入字体}
选择 InkStruct 或 OpenInk 的通用 v1 字库后即可排版。导入旧 InkStruct 字体包时，程序在后台校验并转换，保留原包供旧文档找回原来使用的字形。默认转换包含字形来源和数学轴。需要保留逐点压力和时间时，使用导入按钮下拉菜单中的\ui{导入旧字体并携带压力与时间}。转换失败时，已通过校验的原包仍可使用。

\section{字体库与当前文档（0.1.35）}
\entry{主窗口 → 字体库}
字体库左侧搜索与筛选，右侧分别查看真实笔迹、字符和多份样本、来源与版本、文件信息。缺字显示红框，不用系统字形替代。相同字体身份的版本集中显示，文件较新不代表推荐版本。单击预览与双击详情均不自动更换本文。

\ui{用于当前文档…}先验证资源并进行实际试排，列出缺字、来源和冻结选样保留数量。确认才改变本文，支持撤销重做；检查期间文档或资源变化时拒绝旧结果。核心记录、度量和来源未变的冻结选样保持，其他样本明确重新选择。确认不会解除导出检查。

\ui{设为新文档默认}只影响以后新建的文档。\ui{更多操作}可修改显示名称及移出或恢复列表，不修改字体包和旧文档。导出可选用户版或完整副本，原包保留，不能覆盖已有文件。

\section{较大字体（0.1.46）}
可导入展开总量不超过4 GiB的字体，保留单成员32 MiB、每候选点数、成员数和全部哈希/来源检查。展开总量按ZIP内未压缩数据计算：磁盘文件不大，也可能超过旧版上限；项目保存与排版仍有独立容量保护。

user\_001同字最多8份并包含相同英文、希腊、数学等字符的实测用户字体，全结构库包约611 MB，简体支持字集包约454 MB，压缩体积减少约25.7\%。简体支持集为2013年\ui{通用规范汉字表}与当前结构库的6866字交集；结构库缺字仍列明，不能当作全部简体字已经支持。这是字体体积示例，不是应用安装包体积。

两包展开后都超过旧2 GiB限制，需新版程序读取；升级不自动导入新包或替换旧项目资源。先在字体库选择\ui{用于当前文档…}检查实际影响，确认后再应用。字体能载入不等于来源或数学轴已经人工审阅，原有警告保留。

\section{兼容外部生成目录}
\entry{排版 → 文件 → 高级资源（仅当前文档）}
\begin{enumerate}
\item 点击\ui{选择生成字形目录…}，选择包含 \code{trajectories/} 的生成结果目录，供汉字等内容使用。
\item 在字体库的\ui{制作字体包}中选择同一书写者的生成批次、扩展采集目录或 ZIP，核对清单后按类别入包。
\item 等待预览更新，检查缺字；同一个视觉符号可能对应不同 Unicode 字符，应按诊断补齐实际需要的资源。
\item 保存为 \code{.inkdoc}，保留文档实际使用的资源与选样。
\end{enumerate}

内置单线数字和小数点用于快速试排。它们能完成数字练习，但不表示某位书写者已经录制了这些字符。更换资源后应检查整体笔迹是否一致，不将内置数字误认为个人样本。

\ui{刷新当前外部资源}用于重新读取外部生成目录的变化，程序不会自动监控目录。\ui{载入全部字}用于字库检查，会把生成批次的字符载入检查正文；正式报告操作前先保存，建议另开检查项目使用。

\section{缺字与选样}
不同字符和不同样本各有身份。同一字符可有多个录制写法，排版按样本池选择，并保留已有实例的选择；修改附近正文、重排或重开项目不应让未改内容任意换字。

界面能显示某个字符并不说明存在它的轨迹。红色占位或缺字诊断表示资源不够；不要把占位当作成品。补资源后重新排版；若改的是资源版本，应再次观察实际字形。

对整份文档检查补采需求，可让 Agent 使用附录\ref{ch:agent}的资源检查命令。符号采集 CSV 默认每个缺失字符 5 次，汉字生成需求和不支持字符另列在报告中。空缺口清单仅说明资源存在，不证明根号结构、数学轴或页面尺寸已经合格。

\section{字符校正与旧评价}
在预览中点击字符，右侧\ui{字符校正}可调整位置和缩放。选择应用范围后校正可随文档保存，也可在支持的范围另存校正字体。旧项目中的评价记录继续随项目读取和保存。
